\scanpage{202}
\mspage{L 18}
présenter un tel travail, que ce soient \deleted{\unsure{des}}\inserted{les}
\inserted{manuels et autres} textes didactiques, \deleted{ou manuels,}
ou \deleted{\unsure{des}}\inserted{les} articles \deleted{\unsure{ou}}\inserted{et}
mémoires\deleted{\unsure{,}} \deleted{\unsure{ou de}} originaux, ou les \deleted{exposés} cours oraux
\deleted{ou}\inserted{et} exposés de séminaires etc.
Il y a, depuis des millénaires semblerait-il, depuis les\deleted{ori} origines mêmes de\deleted{\unsure{s}}
la mathématique et des autres sciences, \deleted{que nous cultivons,}
une sorte de ``conspiration du silence'' autour de ces \deleted{aspect}
``\uline{inavouables} \uline{labeurs}'' qui préludent à l'éclosion de toute idée nouvelle,
grande ou petite, venant renouveler notre connaissance d'une portion de\inserted{ce} monde où nous vivons.

Pour tout dire, il semblerait que la répression de la connaissance de cet aspect là\deleted{, de plus}
ou de ce stade-là, le plus crucial de tous\deleted{\unsure{,}} dans le travail de découverte
(et dans le travail créateur en général), soit à tel point efficace, à tel point intériorisé
par ceux-là même qui pourtant connaissent un tel travail de première main, qu'on jurerait souvent
que même ceux-là \deleted{l'ont oublié} en ont éradiqué toute trace de leur souvenir conscient.
Un peu comme dans une société puritaine à outrance, une femme aurait éradiqué de son souvenir,
\deleted{toute trace} en relation à\inserted{chacun de} ces enfants qu'elle se fait un devoir de moucher
et de torcher, le moment de l'étreinte (subie peut-être à contre-coeur) qui le\deleted{\unsure{s}} fit concevoir,
les longs mois de la grossesse\deleted{\unsure{,}} (vécue comme\inserted{une} \unsure{inconvenante})
et les longues heures de l'accouchement (vécues comme un\inserted{peu ragoûtant} martyre,
suivi enfin d'une délivrance).
\margin{\figure{\begin{tikzpicture}[scale=.3]
\draw[blue,->] (0,0)--(2,0)--(2,1.5)--(5,1.5);\draw[blue] (1.8,0)--(1.8,1.5);
\end{tikzpicture}}}
Cette comparaison peut paraître \unsure{outée|outrée}, et elle l'est peut-être,
si je l'applique à ce que je me rappelle à présent de l'esprit que j'ai connu dans le milieu\inserted{mathématique}
dont je faisais moi-même partie, il y a encore vingt ans.
Mais \deleted{dans} j'ai pu me rendre compte, au cours de ma réflexion\deleted{, \unsure{que dans les}}
dans Récoltes et Semailles, et de façon saisissante ces tout derniers mois surtout
(avec l'écriture des Quatre Opérations), qu'il y a eu depuis mon départ de la scène mathématique
une \deleted{dégra} dégradation stupéfiante dans l'esprit qui aujourd'hui fait loi dans les milieux
que j'avais connus, et (me semble-t-il, dans une large mesure au moins) dans le monde mathématique en général.
Il est possible même que mon départ, tant par ma personnalité mathématique très particulière
que par les conditions qui ont entouré ce départ, ait agi comme un catalyseur dans une évolution
qui était déjà en train de se faire ($\unsure{*|x}$); \deleted{\unsure{et}}\inserted{évolution}
dont bien sûr je n'ai rien su percevoir (pas plus qu'aucun de mes collègues et amis,
à la seule exception peut-être de Claude Chevalley).
L'aspect de cette dégradation\inserted{(parmi de nombreux autres ($\unsure{**|xx}$))}
auquel je pense ici est le \deleted{\unsure{dérision ou mépris}}
m\,é\,p\,r\,i\,s\quad t\,a\,c\,i\,t\,e,
\deleted{prenant parfois des formes} quand ce n'est la dérision
\deleted{à peine voilée}\inserted{sans équivoque}, vis à vis de ce qui, dans le travail (mathémati-
\margin{201}
